% Part: first-order-logic
% Chapter: model-theory
% Section: nonstandard-arithmetic

\documentclass[../../../include/open-logic-section]{subfiles}

\begin{document}

\olfileid{mod}{bas}{nsa}
\section{अंकगणिताची अमानक प्रतिमाने}

\begin{defn}
$\Lang{L_N}$ ही अंकगणिताची !!{language} असू द्या. तिच्यात
!!{constant} $\Obj{0}$, द्विस्थानी !!{predicate} $<$,
एकस्थानी !!{function} $\prime$, आणि द्विस्थानी !!{function}s
$+$ व $\times$ आहेत.
\begin{enumerate}
\item अंकगणिताचे \emph{मानक प्रतिमान} म्हणजे $\Lang{L_N}$ साठीची
  !!{structure} $\Struct{N}$: तिचा प्रांत $\Nat = \{ 0, 1, 2, \dots\}$
  असतो; $\Obj{0}$ चा अर्थ $0$, $<$ चा अर्थ $\Nat$ वरील लहान-असण्याचा
  संबंध, तर $\prime$, $+$ आणि $\times$ यांचे अर्थ अनुक्रमे
  $\Nat$ वरील उत्तरवर्ती, बेरीज आणि गुणाकार असे असतात.
\item \emph{सत्य अंकगणित} म्हणजे $\Theory{N}$ ही उपपत्ती.
\end{enumerate}
\end{defn}

$\Lang{L_N}$ मध्ये काम करताना $\Obj{0}$ वर उत्तरवर्ती फलन
$n$ वेळा लावून मिळणाऱ्या $\Obj{0}^{\prime\dots\prime}$
या रूपातील पदाचे संक्षिप्त रूप~$\num{n}$ लिहितो.

\begin{defn}
$\Lang{L_N}$ साठीची !!{structure}~$\Struct{M}$ ही
$\Struct{N} \iso \Struct{M}$ असेल तेव्हा आणि केवळ तेव्हाच
\emph{मानक} असते.
\end{defn}

\begin{thm}\ollabel{thm:non-std}
सत्य अंकगणिताची अमानक !!{enumerable} प्रतिमाने आहेत.
\end{thm}

\begin{proof}
$\Lang{L_N}$ मध्ये नवा !!{constant}~$c$ घालून तिचा विस्तार करा
आणि पुढील उपपत्ती विचारात घ्या:
\[
\Theory{N} \cup \Setabs{\num{n} < c}{n \in \Nat}.
\]
ही उपपत्ती सांत पूर्ततायोग्य आहे. त्यामुळे संहततेनुसार तिला
$\Struct{M}$ हे प्रतिमान आहे; अधोगामी लोव्हेनहाइम--स्कोलेम
प्रमेयानुसार ते !!{enumerable} घेता येते. $\Domain{M}$ हा
$\Struct{M}$ चा प्रांत असताना, $\Lang{L}$ मधील अतार्किक
स्थिरांकांचे $\Struct{M}$ मधील अर्थनिर्धारण असे लिहू:
$\mathbf{z} = \Assign{\Obj{0}}{M} \in \Domain{M}$,
${\prec} = \Assign{<}{M} \subseteq M^2$,
$* = \Assign{\prime}{M} \colon \Domain{M} \to \Domain{M}$,
आणि $\oplus = \Assign{+}{M}, \otimes =
\Assign{\times}{M} \colon \Domain{M}^2 \to \Domain{M}$.
प्रत्येक $x \in \Domain{M}$ साठी, $x$ वर~$*$ लावल्याने
मिळणाऱ्या $\Domain{M}$ मधील घटकाला $x^*$ लिहू.
%READERNOTE{OLINC-334: स्रोत येथे आधी निश्चित केलेल्या L_N ऐवजी L वापरतो आणि संबंधाचा प्रांत M^2 असा लिहितो; ही संकेतविसंगती जशीच्या तशी ठेवली आहे.}

आता, $h$ हे $\Struct{N}$ आणि $\Struct{M}$ यांचे
समरूपण असेल, तर काही $n \in \Nat$ साठी
$h(n) = \Assign{c}{M}$ असते. म्हणून $\Struct{N}$ मध्ये
$s(x) = n$ अशी कोणतीही चलनियुक्ती $s$ घ्या. मग
$\Sat{N}{\eq[\num{n}][x]}[s]$; तसेच
\olref[iso]{thm:isom} च्या सिद्धतेनुसार
$\Sat{M}{\eq[\num{n}][x]}[h \circ s]$.
त्यामुळे $\Assign{c}{M} = \mathbf{z}^{*\cdots *}$
(येथे $*$ ही क्रिया $n$ वेळा केली आहे).
पण गृहीतकानुसार $\Sat{M}{\num{n} < c}$ आहे आणि
$\prec$ हा संबंध अपरावर्ती आहे; म्हणून हे अशक्य आहे.
त्यामुळे $\Struct{M}$ अमानक आहे.
\end{proof}

\begin{prob}
संच $X$ वरील संबंध $R$ हा तेव्हा आणि केवळ तेव्हाच
\emph{सुस्थापित} असतो, जेव्हा $R$ मध्ये अनंत अवरोही साखळी
नसते; म्हणजे $\dots x_2Rx_1Rx_0$ अशी अट पूर्ण करणारे
$X$ मधील $x_0$, $x_1$, $x_2$,~\dots नसतात.
झेर्मेलो--फ्रेंकेल संच उपपत्ती~$ZF$ सुसंगत आहे असे मानून,
$ZF$ ची असुस्थापित प्रतिमाने आहेत हे दाखवा; म्हणजे
$\dots x_2 \in x_1 \in x_0$ असणारी प्रतिमाने
$\mathfrak{M}$ आहेत हे सिद्ध करा.
\end{prob}

\olref{thm:non-std} मधील अमानक प्रतिमान
$\Struct{M}$ हे मानक प्रतिमानाशी प्राथमिक सममूल्य असल्यामुळे
$\Struct{M}$ चे अनेक गुणधर्म काढता येतात. या विभागाचा
उर्वरित भाग अशा गुणधर्मांच्या अभ्यासासाठी आहे; त्यातून
$\Theory{N}$ च्या अमानक !!{enumerable} प्रतिमानांचे
अचूक लक्षण मिळेल.

\begin{enumerate}
\item $\Domain{M}$ मधील कोणताही घटक स्वतःहून $\prec$-लहान
  नसतो: $\lforall[x][\lnot x < x]$ हे वाक्य $\Struct{N}$
  मध्ये सत्य असल्यामुळे $\Struct{M}$ मध्येही सत्य आहे.
\item याच रीतीने $\prec$ हा $\Domain{M}$ वरील
  \emph{रेषीय क्रम} आहे, म्हणजे तो $\Domain{M}$ वरील
  पूर्ण, अपरावर्ती आणि संक्रमक संबंध आहे.
\item $\mathbf{z}$ हा $\Domain{M}$ मधील $\prec$-लघुतम घटक आहे.
\item $\Domain{M}$ मधील प्रत्येक घटक त्याच्या $*$-उत्तरवर्तीपेक्षा
  $\prec$-लहान असतो; आणि $x^*$ हा $x$ पेक्षा मोठ्या
  $\Domain{M}$ मधील घटकांपैकी $\prec$-लघुतम असतो.
\item $\Struct{M}$ मध्ये $\Nat$ शी समरूप असा $\prec$ चा
  आरंभीचा खंड आहे: $\mathbf{z}, \mathbf{z}^*,
  \mathbf{z}^{**}, \dots$. त्याला $\Domain{M}$ चा
  \emph{मानक भाग} म्हणतो. $\Domain{M}$ मधील इतर प्रत्येक
  घटक \emph{अमानक} असतो. $\Domain{M}$ मध्ये अमानक
  घटक असलेच पाहिजेत;
  अन्यथा \olref{thm:non-std} च्या सिद्धतेतील $h$ हे
  समरूपण ठरेल. $n, m, \dots$ ही !!{variable}s आपण
  $\Struct{M}$ च्या या मानक भागावर चालणारी मानतो.
\item प्रत्येक अमानक घटक प्रत्येक मानक घटकापेक्षा मोठा
  असतो; कारण प्रत्येक $n \in \Nat$ साठी
  \[
  \Sat{N}{\lforall[z][(\lnot(\eq[z][\Obj{0}] \lor
  \dots \lor \eq[z][\num{n}]) \lif \num{n} < z)]},
  \]
  म्हणून $z \in \Domain{M}$ सर्व मानक घटकांहून
  वेगळा असेल, तर तो त्या सर्वांपेक्षा \emph{मोठा} असतो.
\item $\mathbf{z}$ सोडून $\Domain{M}$ मधील प्रत्येक
  घटक हा एखाद्या एकमेव $\Domain{M}$ मधील घटकाचा
  $*$-उत्तरवर्ती असतो; त्या पूर्ववर्तीला $^*x$ लिहितो.
  $x = y^*$ असेल, तर $x$ आणि $y$ यांपैकी एक मानक
  असल्यास दोन्ही मानक असतात (आणि एक अमानक असल्यास
  दोन्ही अमानक असतात).
\item $\Domain{M}$ वर $\approx$ हा सममूल्यता संबंध
  परिभाषित करा: एखाद्या \emph{मानक} $n$ साठी
  $x \oplus n = y$ किंवा $y \oplus n =x$
  असेल तेव्हा आणि केवळ तेव्हाच $x \approx y$.
  दुसऱ्या शब्दांत, $x$ आणि $y$ यांच्यात सांत अंतर
  असेल तेव्हा आणि केवळ तेव्हाच $x \approx y$.
  $n$ आणि $m$ मानक असतील, तर $n \approx m$.
  $x$ चा \emph{खंड} म्हणजे सममूल्यता वर्ग
  $[x] = \Setabs{y}{x \approx y}$.
\item $x \prec y$ पण $x \not\approx y$ असे समजा.
  $\Sat{N}{\lforall[x][\lforall[y][(x < y \lif (x' < y \lor x' =
      y))]]}$ असल्यामुळे, $x^* \prec y$ किंवा
  $x^* = y$. दुसरी शक्यता $x \approx y$ सूचित
  करत असल्यामुळे अशक्य आहे; म्हणून
  $x \prec y$. त्याचप्रमाणे $x \prec y$ आणि
  $x \not\approx y$ असतील, तर $x \prec {^*y}$.
  म्हणून $x \prec y$ आणि $x \not\approx y$ असतील,
  तर प्रत्येक $w \approx x$ हा प्रत्येक $v \approx y$
  पेक्षा $\prec$-लहान असतो. त्यानुसार, प्रत्येक
  अमानक खंड $[x]$ पुढील द्विदिश अनंत साखळी बनवतो:
  \[
  \cdots \prec  {^{**}x} \prec {^*}x \prec x \prec x^* \prec x^{**}
  \prec \cdots
  \]
  पूर्णांकांच्या क्रमप्रकारासारखा असल्यामुळे तिला
  $Z$-साखळी म्हणतात.
  %READERNOTE{OLINC-336: गोठवलेल्या स्रोताचा 'प्रत्येक खंड' हा दावा मानक खंडासाठी चुकीचा आहे; द्विदिश अनंत साखळी फक्त अमानक खंडांना लागू होते.}
  %READERNOTE{OLINC-337: स्रोताचे 'म्हणून x<y' हे वाक्य आधीचेच गृहीतक पुन्हा सांगते; त्या ठिकाणी मिळालेला अधिक मजबूत निष्कर्ष x चा उत्तरवर्ती y पेक्षा लहान आहे.}
\item $\prec$ हा क्रम खंडांवरही नेऊ शकतो: $x \prec y$
  आणि ते भिन्न खंडांत असतील, तर $x$ चा खंड
  $y$ च्या खंडापेक्षा लहान असतो. अमानक घटक असणारा
  खंड \emph{अमानक} असतो. मानक खंड हा लघुतम खंड आहे.
  %READERNOTE{OLINC-338: भिन्न खंडांसाठी x आणि y असण्याची अट स्रोताच्या पहिल्या वाक्यात गाळलेली आहे; एकाच खंडातील x<y त्यांचा खंड स्वतःहून लहान करत नाही.}
\item लघुतम अमानक खंड नाही: $y$ अमानक असेल,
  तर $x \prec y$ आणि $x \not\approx y$ अशी अट
  पूर्ण करणारा अमानक $x$ असतो. सिद्धता: मानक
  प्रतिमान $\Struct{N}$ मध्ये प्रत्येक संख्येला दोनने
  भाग जातो, कदाचित बाकी एक उरते; म्हणजे
  $\Sat{N}{\lforall[y][\lexists[x][(y = x + x \lor y = x + x +
        \Obj{0}')]]}$. प्राथमिक सममूल्यतेमुळे
  प्रत्येक $y \in \Domain{M}$ साठी असा
  $x \in \Domain{M}$ असतो की $x \oplus x = y$
  किंवा $x \oplus x \oplus \mathbf{z}^*= y$.
  $x$ मानक असेल, तर $y$ सुद्धा मानक असेल;
  म्हणून $x$ अमानक आहे. शिवाय $x$ आणि $y$
  वेगळ्या खंडांत येतात, म्हणजे $x \not\approx y$.
  तसे नसल्यास, एखाद्या मानक $n$ साठी
  $x \oplus n = y$ असेल. $y = x \oplus x$
  असेल, तर $x \oplus n = x \oplus x$; बेरीज
  रद्द करण्याच्या नियमानुसार $x = n$ येईल.
  हा नियम $\Struct{N}$ मध्ये आणि म्हणून
  $\Struct{M}$ मध्येही चालतो; पण मग $x$
  मानक ठरेल. $y = x \oplus x \oplus \mathbf{z}^*$
  या प्रसंगातही तसेच होते.
  %READERNOTE{OLINC-335: स्रोताच्या अर्धे करण्याच्या सूत्रात प्रत्येक y साठी 'प्रत्येक x' असे चुकीने आहे; अस्तित्वदर्शक x आवश्यक आहे आणि येथे दुरुस्त केला आहे.}
\item याच प्रकारच्या युक्तिवादाने सर्वांत मोठा खंडही नाही.
\item खंडांचा क्रम सघन आहे: $[x]$ हा $[y]$
  पेक्षा लहान असेल ($x \not\approx y$), तर
  त्यांच्या मध्ये दोन्हींपेक्षा वेगळा खंड $[z]$
  असतो. $x \prec y$ असे समजा. आधीप्रमाणे
  $x \oplus y$ ला दोनने भाग जातो (कदाचित
  बाकी उरते); म्हणून काही $u \in \Domain{M}$
  साठी $x \oplus y = u \oplus u$ किंवा
  $x \oplus y = u \oplus u \oplus \mathbf{z}^*$.
  $u$ हा $x$ आणि $y$ यांचा मध्य आहे आणि
  त्यांच्यामध्ये येतो. $x \oplus y = u \oplus u$
  असे समजा (दुसरा प्रसंगही तसाच): $u \approx x$
  असेल, तर एखाद्या मानक $n$ साठी
  \[
  x \oplus y = x \oplus n \oplus x \oplus n,
  \]
  म्हणून $y = x \oplus n \oplus n$; यातून
  $x \approx y$ येते, जे गृहीतकाविरुद्ध आहे.
  म्हणून $u \not\approx x$. त्याच युक्तिवादाने
  $u \not\approx y$.
\end{enumerate}
त्यामुळे अमानक खंडांचा क्रम परिमेय संख्यांप्रमाणे
आहे: तो अंत्यबिंदू नसलेला !!a{enumerable} रेषीय क्रम आहे.
यावरून सत्य अंकगणिताच्या कोणत्याही दोन
!!{enumerable} अमानक प्रतिमानांचे,
$\Struct{M}_1$ आणि $\Struct{M_2}$ यांचे,
फक्त $<$ आणि $=$ असणाऱ्या भाषेतील
न्यूनीकृत प्रतिमान समरूप असतात.
खरोखर, समरूपण $h$ असे बनवता येते:
$\Struct{M_1}$ आणि $\Struct{M_2}$ यांचे
मानक भाग मानक प्रतिमान $\Struct{N}$ शी,
म्हणून एकमेकांशी, समरूप असतात. अमानक भाग
बनवणारे खंड परिमेयांप्रमाणे क्रमबद्ध असतात;
म्हणून \olref[dlo]{thm:cantorQ} नुसार ते
समरूप असतात. प्रत्येक खंडातील कोणतेही
घटक जुळवून खंडांचे समरूपण खंडांच्या
\emph{आत} विस्तारता येते; मग
$\Struct{M_1}$ मधील $x$ च्या उत्तरवर्तीची
प्रतिमा ही $\Struct{M_2}$ मधील $x$ च्या
प्रतिमेचा उत्तरवर्ती ठरवता येते.
यावरून मात्र $\mathfrak{M}_1$ आणि
$\mathfrak{M}_2$ ही संपूर्ण अंकगणितीय
भाषेत समरूप आहेत असे \emph{निघत नाही}
(कारण समरूपता नेहमी चिन्हसंचाच्या संदर्भातच
असते). $\Domain{M_1}$ आणि $\Domain{M_2}$
वर बेरीज व गुणाकार परिभाषित करण्याचे
असमरूप मार्ग आहेत. एका संपूर्ण उपपत्तीच्या
गणनीय प्रतिमानांची संख्या नेमकी~2
असू शकत नाही, या वॉट यांच्या प्रसिद्ध
प्रमेयावरूनही हे निघते.

\begin{prob}
अंकगणिताच्या अमानक प्रतिमानात सर्वांत मोठा
खंड असू शकत नाही हे दाखवा.
\end{prob}

\begin{prob}
$\Lang{L}$ ही $<$ हा एकच !!{predicate}
($\eq$ व्यतिरिक्त) असणारी प्रथम-क्रमाची
!!{language} असू द्या आणि
$\Struct{N} = (\Nat, <)$ असू द्या.
$\Struct{N}$ चे सर्व सांत किंवा सहसांत
उपसंच परिभाषणीय आहेत. \emph{हेच एकमेव}
$\Struct{N}$ चे परिभाषणीय उपसंच आहेत हे दाखवा.

(सूचना: प्रथम, $\Obj{prc}(x,y)$ हे
$\Lang{L}$ मधील “$x$ हा $y$ चा लगतचा
पूर्ववर्ती आहे” याचे संक्षिप्त रूप असलेले
सूत्र असू द्या:
\[
x<y \land \lnot \lexists[z][(x<z \land z < y)].
\]
आता $\Struct{N}$ च्या कोणत्याही परिभाषणीय
उपसंचाला $\Lang{L}$ मधील सूत्र $!A(x)$
अनुरूप असते. अशा प्रत्येक $!A$ साठी
पुढील वाक्य $!D$ विचारात घ्या:
\[
\lexists[x][\lforall[y][\lforall[z][((x<y \land x<z \land \Obj{prc}(y,z)
  \land !A(y)) \lif !A(z))]]].
\]
$\Sat{N}{!D}$ असेल तेव्हा आणि केवळ
तेव्हाच $!A$ ने परिभाषित केलेला
$\Struct{N}$ चा उपसंच सांत किंवा
सहसांत असतो हे दाखवा.

आता $\Struct{M}$ हे $\Struct{N}$
शी प्राथमिक सममूल्य अमानक प्रतिमान असू द्या.
$a \in \Domain{M}$ अमानक असेल, तर
$b, c \in \Domain{M}$ हे $a$ पेक्षा
मोठे घ्या आणि $b$ हा~$c$ चा लगतचा
पूर्ववर्ती असू द्या. मग $\Domain{M}$
चे असे स्वसमरूपण $h$ आहे की $h(b)=c$
(का?). त्यामुळे $b$ ने $!A$ ची पूर्ति
केली, तर $c$ नेही ती केली पाहिजे (का?).
म्हणून $!D$ हे $\Struct{M}$ मध्ये सत्य
आहे आणि त्यामुळे $\Struct{N}$ मध्येही.
पण याचा अर्थ असा की $!A$ ने परिभाषित
केलेला $\Struct{N}$ चा उपसंच सांत
किंवा सहसांत आहे.
\end{prob}

\end{document}
